\section{Introduction}
\label{sec:introduction}

The minimum weighted feedback arc set (MWFAS) problem asks for a minimum-weight set of arcs whose removal makes a directed graph acyclic; equivalently, it asks for a vertex ordering with minimum total backward-arc weight. The problem is NP-hard \cite{K72} and arises whenever directed pairwise relations contain inconsistent cycles, including precedence aggregation, dependency repair, scheduling, and ranking from directed evidence \cite{F90,ACN08}. In weighted settings, not all violations are interchangeable: reversing a high-weight precedence relation is more costly than reversing a weak one, so useful heuristics must react to both graph structure and arc-weight asymmetry.

This paper focuses on \emph{sparse nonnegative weighted digraphs}. That regime differs structurally from dense linear-ordering or tournament instances. In many sparse graphs, large regions are already acyclic or only weakly cyclic, while the hard part of the objective is concentrated in a smaller set of strongly connected components (SCCs). Dense complete ordering instances behave differently: every ordered pair carries weight, the cyclic core effectively spans the whole graph, and matrix-based pairwise-ordering heuristics are naturally aligned with the input model. A method that is strong on sparse digraphs therefore need not be strongest on dense ordering benchmarks.

This creates a practical gap. Exact formulations provide certification but do not scale broadly enough to serve as the main tool on the sparse benchmark studied here \cite{BSNA21}. Classical greedy heuristics remain fast and interpretable, yet they may spend effort too globally or may not exploit the fact that residual backward weight is often localized after a strong constructive pass. Earlier local-ratio and weighted minimal/stable constructions provide useful seeds, but they leave open whether a \emph{protected} SCC-local refinement stage can improve solutions consistently without ever returning a worse answer than the best available seed.

The central contribution is \emph{incumbent-protected SCC neighborhood search} (IPSNS). IPSNS begins from the better of two attributed seeds, scores nontrivial SCCs by their current backward-weight contribution, and applies SCC-local destroy-and-repair moves inside fixed original-graph neighborhoods. The destroy phase reactivates heavy removed internal arcs and removes light active internal arcs. The repair phase reruns local-ratio reduction and heavy-first add-back only inside the selected SCC neighborhood. A candidate is accepted only if the repaired global ordering is feasible and strictly improves backward weight; otherwise the incumbent is restored exactly. IPSNS is therefore stochastic in neighborhood choice but monotone in returned objective value.

The two seeds are supporting components rather than coequal contributions. LR-TA (local-ratio topological add-back) is an engineered sparse-graph seed derived from the directed local-ratio lineage of Demetrescu and Finocchi \cite{DF03} and connected to earlier author work on ranking-oriented MWFAS \cite{VahidiKoutis2024arxiv}. The WMSF-style seed is an adapted implementation of the weighted minimal/stable feedback arc set line of Cavallaro and Cutello \cite{CC25}. The present paper does not claim those foundations as new; it uses them as complementary entry points whose better incumbent IPSNS then protects and selectively improves.

The empirical message is intentionally narrow. On the primary sparse comparison, IPSNS attains the minimum observed backward weight among the evaluated methods on 96 of 97 standard nonnegative instances. On the 57-instance exact-validation subset, it matches the dynamic-programming optimum on 56 instances. On the 93-instance common sparse subset, repeated-run per-instance medians over 20 runs yield 38 wins, 55 ties, and 0 losses for IPSNS against the evaluated DRMacIver/FAS executable. A dense LOLIB transfer study marks the regime boundary: matrix-based ordering methods remain preferable on complete weighted ordering instances. The paper's claim is therefore that protected SCC-local refinement is an effective sparse-graph algorithm-engineering strategy among the evaluated methods, not that it solves every MWFAS regime equally well.

The contributions are:
\begin{itemize}
    \item \textbf{Algorithmic contribution.} We introduce IPSNS, an incumbent-protected SCC-local destroy-and-repair heuristic for sparse nonnegative weighted digraphs. The contribution lies in the problem-specific integration of fixed SCC neighborhoods, contribution-based SCC selection, two-sided perturbation, SCC-restricted repair, and strict global-improvement acceptance.
    \item \textbf{Integrated framework.} We evaluate IPSNS together with two explicitly attributed seeds: an engineered LR-TA seed derived from prior local-ratio work and a WMSF-style seed adapted from Cavallaro and Cutello's weighted minimal/stable pipeline.
    \item \textbf{Supporting analysis.} We formalize implementation-level properties for seed feasibility, one-pass add-back minimality, and IPSNS incumbent non-worsening, together with concise complexity summaries. These are supporting guarantees for the implemented framework rather than new approximation results.
    \item \textbf{Computational evidence.} We report sparse benchmark comparisons, exact and MIP validation, ablation and holdout-supported parameter checks, repeated-run robustness against an external executable with run-to-run variation, and a dense LOLIB scope-boundary study.
\end{itemize}

The remainder of the paper positions the work, defines the objective, presents the framework, summarizes the experimental design, reports the claim-bearing results, discusses the sparse and dense regimes, and concludes.
